\chapter{量子力学核心计算：边界匹配与近似展开}

\section{概率流守恒}
对
\[
i\hbar\partial_t\psi=
-\frac{\hbar^2}{2m}\nabla^2\psi+V\psi
\]
取复共轭（$V$ 实）：
\[
-i\hbar\partial_t\psi^*
=-\frac{\hbar^2}{2m}\nabla^2\psi^*+V\psi^*.
\]
第一式乘 $\psi^*$，第二式乘 $\psi$ 后相减：
\[
i\hbar(\psi^*\partial_t\psi+\psi\partial_t\psi^*)
=-\frac{\hbar^2}{2m}
(\psi^*\nabla^2\psi-\psi\nabla^2\psi^*).
\]
左边为 $i\hbar\partial_t|\psi|^2$。利用
\[
\psi^*\nabla^2\psi-\psi\nabla^2\psi^*
=\nabla\cdot(\psi^*\nabla\psi-\psi\nabla\psi^*)
\]
得到
\[
\boxed{\partial_t\rho+\nabla\cdot\vb j=0},
\quad
\rho=|\psi|^2,\quad
\vb j=\frac{\hbar}{2mi}
(\psi^*\nabla\psi-\psi\nabla\psi^*).
\]
积分并要求无穷远通量为零，得总概率守恒。

\section{有限势垒的边界匹配}
势垒 $V=V_0$ 位于 $0<x<a$，$E<V_0$。三区波函数
\[
\psi_I=e^{ikx}+re^{-ikx},\quad
\psi_{II}=Ae^{\kappa x}+Be^{-\kappa x},\quad
\psi_{III}=te^{ikx},
\]
\[
k=\sqrt{2mE}/\hbar,\qquad
\kappa=\sqrt{2m(V_0-E)}/\hbar.
\]
有限阶跃处积分 Schrödinger 方程可知 $\psi,\psi'$ 连续。于
$x=0,a$ 写四个线性方程并消去 $A,B,r$，得到
\[
\boxed{
T=|t|^2=
\left[1+\frac{V_0^2\sinh^2(\kappa a)}
{4E(V_0-E)}\right]^{-1}}.
\]
厚势垒 $\kappa a\gg1$：
\[
T\simeq
\frac{16E(V_0-E)}{V_0^2}e^{-2\kappa a}.
\]
指数来自禁阻区振幅衰减，前因子来自两界面匹配。

\section{谐振子升降算符}
定义
\[
a=\sqrt{\frac{m\omega}{2\hbar}}x+
\frac{i}{\sqrt{2m\hbar\omega}}p,
\quad
a^\dagger=\sqrt{\frac{m\omega}{2\hbar}}x-
\frac{i}{\sqrt{2m\hbar\omega}}p.
\]
用 $[x,p]=i\hbar$：
\[
[a,a^\dagger]
=-2i\sqrt{\frac{m\omega}{2\hbar}}
\frac1{\sqrt{2m\hbar\omega}}[x,p]=1.
\]
反解 $x,p$ 代回 Hamiltonian：
\[
H=\hbar\omega\left(a^\dagger a+\frac12\right).
\]
令 $N=a^\dagger a$。由
$[N,a^\dagger]=a^\dagger$，
$a^\dagger|n\rangle$ 的本征值为 $n+1$。又
\[
\|a|n\rangle\|^2
=\langle n|a^\dagger a|n\rangle=n\ge0.
\]
下降链必须在 $n=0$ 终止，否则会出现负范数，故
\[
E_n=\hbar\omega(n+\tfrac12).
\]
基态由 $a\psi_0=0$：
\[
\left(\dv{}{x}+\frac{m\omega}{\hbar}x\right)\psi_0=0
\Rightarrow
\psi_0\propto e^{-m\omega x^2/2\hbar}.
\]

\section{角动量量子化}
\[
[J_i,J_j]=i\hbar\epsilon_{ijk}J_k,\quad
J_\pm=J_x\pm iJ_y.
\]
可得
\[
[J_z,J_\pm]=\pm\hbar J_\pm,\qquad
[J^2,J_\pm]=0.
\]
故 $J_\pm|jm\rangle$ 保持 $j$、把 $m$ 改变 $\pm1$。范数
\[
\|J_+|jm\rangle\|^2
=\langle jm|J_-J_+|jm\rangle
=\hbar^2[j(j+1)-m(m+1)].
\]
最高态要求范数为零，得 $m_{\max}=j$；最低态得
$m_{\min}=-j$。从 $j$ 到 $-j$ 的步数 $2j$ 必为整数，所以
$j=0,\frac12,1,\ldots$。

\section{非简并微扰逐阶展开}
\[
H=H_0+\lambda V,\quad
|n\rangle=|n^{(0)}\rangle+\lambda|n^{(1)}\rangle+\cdots,
\]
\[
E_n=E_n^{(0)}+\lambda E_n^{(1)}
+\lambda^2E_n^{(2)}+\cdots.
\]
代入 $(H_0+\lambda V)|n\rangle=E_n|n\rangle$。一阶：
\[
(H_0-E_n^{(0)})|n^{(1)}\rangle
=-(V-E_n^{(1)})|n^{(0)}\rangle.
\]
左乘 $\langle n^{(0)}|$，左边为零：
\[
\boxed{E_n^{(1)}=\langle n^{(0)}|V|n^{(0)}\rangle}.
\]
左乘 $\langle m^{(0)}|,\ m\ne n$：
\[
(E_m^{(0)}-E_n^{(0)})
\langle m|n^{(1)}\rangle=-V_{mn},
\]
\[
|n^{(1)}\rangle=\sum_{m\ne n}
\frac{V_{mn}}{E_n^{(0)}-E_m^{(0)}}|m^{(0)}\rangle
\]
（采用中间归一化 $\langle n^{(0)}|n\rangle=1$）。
二阶方程左乘 $\langle n^{(0)}|$：
\[
\boxed{
E_n^{(2)}=\sum_{m\ne n}
\frac{|V_{mn}|^2}{E_n^{(0)}-E_m^{(0)}}}.
\]
分母显示近简并时展开失效。

\section{简并微扰的投影推导}
令 $P$ 投影到简并能量 $E_0$ 的子空间，$Q=1-P$。一阶方程投影：
\[
PVP|\psi^{(0)}\rangle
=E^{(1)}|\psi^{(0)}\rangle.
\]
因此必须先在简并子空间对角化 $PVP$。二阶消去 $Q$ 子空间可写
\[
H_{\rm eff}^{(2)}
=-PVQ\frac1{Q(H_0-E_0)Q}QVP.
\]
Hubbard 强耦合的超交换正是这个公式的应用。

\section{含时微扰与黄金规则}
相互作用绘景
\[
i\hbar\pdv{|\psi_I\rangle}{t}=V_I(t)|\psi_I\rangle.
\]
积分方程
\[
|\psi_I(t)\rangle=|\psi(0)\rangle
-\frac{i}{\hbar}\int_0^tV_I(t')|\psi_I(t')\rangle\dd t'.
\]
第一次迭代给一阶振幅
\[
c_f^{(1)}(t)=-\frac{i}{\hbar}
\int_0^t\langle f|V(t')|i\rangle
e^{i\omega_{fi}t'}\dd t'.
\]
若 $V_{fi}$ 恒定，
\[
|c_f|^2=\frac{|V_{fi}|^2}{\hbar^2}
\frac{4\sin^2(\omega_{fi}t/2)}{\omega_{fi}^2}.
\]
长时间函数峰高 $\sim t^2$、宽 $\sim1/t$、面积 $\sim t$，分布极限
\[
\frac{4\sin^2(\omega t/2)}{\omega^2}
\longrightarrow2\pi t\,\delta(\omega).
\]
对连续末态乘态密度并除以 $t$：
\[
\boxed{\Gamma_{i\to f}
=\frac{2\pi}{\hbar}|V_{fi}|^2\rho(E_f)}.
\]

\section{WKB 展开}
定态方程中令
\[
\psi(x)=\exp\left[\frac{i}{\hbar}S(x)\right].
\]
代入并除以 $\psi$：
\[
\frac{1}{2m}(S')^2-\frac{i\hbar}{2m}S''+V=E.
\]
展开 $S=S_0+(\hbar/i)S_1+\cdots$。零阶：
\[
S_0'=\pm p(x),\qquad p=\sqrt{2m(E-V)}.
\]
一阶：
\[
2S_0'S_1'+S_0''=0
\Rightarrow S_1=-\frac12\ln p.
\]
故允许区
\[
\psi\simeq\frac1{\sqrt p}
\left(C_+e^{i\int p\dd x/\hbar}
+C_-e^{-i\int p\dd x/\hbar}\right).
\]
两转向点连接相位各贡献 $\pi/4$，得
\[
\int_{x_1}^{x_2}p(x)\dd x
=\left(n+\frac12\right)\pi\hbar.
\]

\section{二次量子化推导单体算符}
场算符展开
\[
\hat\psi(\vb r)=\sum_\alpha\phi_\alpha(\vb r)a_\alpha.
\]
单粒子 Hamiltonian 的多体形式
\[
\hat H_1=\int\dd^3r\,
\hat\psi^\dagger(\vb r)h(\vb r)\hat\psi(\vb r).
\]
代入展开：
\[
\hat H_1=\sum_{\alpha\beta}
a_\alpha^\dagger a_\beta
\int\phi_\alpha^*h\phi_\beta\,\dd^3r
=\sum_{\alpha\beta}h_{\alpha\beta}
a_\alpha^\dagger a_\beta.
\]
若 $\phi_\alpha$ 是 $h$ 本征态，$h_{\alpha\beta}=\epsilon_\alpha\delta_{\alpha\beta}$，
\[
\hat H_1=\sum_\alpha\epsilon_\alpha n_\alpha.
\]

\begin{selftestbox}
\begin{enumerate}
\item 概率流推导为何要求势为实数？
\item 势垒边界为何要求 $\psi'$ 连续？
\item 谐振子能谱离散性如何来自范数非负？
\item 角动量梯子为什么必须终止？
\item 微扰论中采用了什么归一化约定？
\item 简并微扰为何等价于求 $PVP$ 本征态？
\item 黄金规则的 $\delta$ 函数来自哪个极限？
\item WKB 振幅为何是 $1/\sqrt p$？
\item 二次量子化如何自动编码交换对称性？
\item Hubbard 超交换与二阶投影公式怎样对应？
\end{enumerate}
\end{selftestbox}
